\section{Cumplimiento normativo, auditoría y puesta en producción}
\subsection{Logging y trazabilidad}
Cada policy update debe acompañarse de un registro: seed, hyper parameters, KL divergence, critic loss. Las conclusiones deben vincularse a una `audit note` para que el regulator documente las assessments.

\begin{table}[H]
\centering
\small
\begin{tabular}{p{0.3\textwidth}p{0.6\textwidth}}
\toprule
Artifact & Contenido \\
\midrule
Prompt registry & Registra el prompt template + temperature \\
Override log & Registra el clinician override + su justificación \\
Evidence patch bank & Almacena el attention map y los highlight frames \\
Release ticket & Incluye el rollout plan + rollback criteria \\
\bottomrule
\end{tabular}
\caption{Audit artifacts del policy gradient release}
\end{table}

\subsection{Checklist de entrega de ingeniería}
Antes de cada release debe completarse:
\begin{enumerate}
    \item Congelar el prompt + hyper parameters;
    \item Verificar que KL divergence < threshold;
    \item Archivar la reward curve + entropy log;
    \item Registrar la fail-safe policy + rollback plan.
\end{enumerate}

\begin{importantbox}{El valor del checklist de release}
\begin{itemize}
    \item Garantiza que cada update tenga una ruta de decisiones trazable;
    \item Acelera el proceso de regulator review;
    \item Forma una knowledge base de mejora continua.
\end{itemize}
\end{importantbox}

\subsection{Resumen de la sección}
El cumplimiento normativo y la auditoría son el último tramo de la puesta en producción del policy gradient: los registros, los artifacts y el checklist conforman el evidence trail aceptable para el regulator.
